\documentclass{cumcmthesis}
% 电子版提交时改为：
% \documentclass[withoutpreface,bwprint]{cumcmthesis}

\title{论文题目}
\tihao{A}                       % 题号：A/B/C/D/E
\baominghao{000000000000}       % 12 位报名参赛队号
\schoolname{学校全称}
\membera{队员一}
\memberb{队员二}
\memberc{队员三}
\supervisor{指导教师}
\yearinput{2026}
\monthinput{09}
\dayinput{13}

\begin{document}

\maketitle

\begin{abstract}
在此填写摘要。摘要内容应包含研究对象、主要方法、关键结果与结论，原则上不超过一页。

\keywords{关键词一\quad 关键词二\quad 关键词三}
\end{abstract}

% 2026 年规范要求正文不要目录，请勿启用 \tableofcontents。

\section{问题重述}
在此准确、简洁地重述赛题。

\section{问题分析}
在此说明研究思路、问题分解和总体技术路线。

\section{模型假设}
\begin{enumerate}
    \item 假设一；
    \item 假设二。
\end{enumerate}

\section{符号说明}
\begin{table}[H]
    \centering
    \caption{主要符号说明}
    \begin{tabular}{cc}
        \toprule[0.8pt]
        符号 & 含义 \\
        \midrule[0.5pt]
        $x$ & 示例变量 \\
        \bottomrule[0.8pt]
    \end{tabular}
\end{table}

\section{模型建立与求解}
在此给出模型、算法、求解过程和结果。

\section{模型检验与分析}
在此给出误差分析、敏感性分析、稳健性检验或其他验证。

\section{模型评价与推广}
在此说明模型的优点、局限与可推广方向。

\begin{thebibliography}{9}
    \bibitem{ref1}
    作者.
    \newblock 文献题目.
    \newblock 出版信息, 年份.
\end{thebibliography}

\section*{人工智能工具使用声明}
请依据竞赛期间有效的全国组委会规定填写；未使用时也应按规定作出明确声明。

\begin{appendices}

\section{支撑材料文件列表}
请逐项列出文件名、格式及用途。若没有支撑材料，请按官方规范明确说明“本论文没有支撑材料”。

\section{源程序}
请附上全部完整、可运行的源程序。若没有使用程序，请按官方规范明确说明“本论文没有用到程序”。

\end{appendices}

\end{document}
